\section{Trust assumptions and remaining work}
\label{sec:boundaries}

The execution result depends on the Lean kernel accepting the proof, the stated Lean axioms, and the adequacy of the formal decoder, verifier profile, integer and floating-point semantics, and Talos execution model for the claimed WebAssembly behavior.  The proof connects those definitions to the embedded bytes.  Native deployment further depends on the external-file check, runtime execution, host byte transfers and calls, operating system, and hardware.  Arithmetic NaN canonicalization is enabled in the retained Wasmtime configuration.  The package's empty host-import assumption list concerns imported WebAssembly functions.  It does not remove the external runtime and orchestration assumptions.

The numerical result adds assumptions about the input checkpoint relation, arithmetic ranges, positive denominators, magnitudes, source-intermediate identities, successful steps, and native checker evaluation.  Source-level error theorems enter the quantized package's declaration audit alongside the execution theorems.  The principal exact-binary theorem in Equation~\eqref{eq:artifact} states execution behavior.  Applying the numerical comparison to an observed deployment also requires the numerical premises and the connection from measured operands to the specified source computation.

The experiments compare one pinned model on a small fixed collection of inputs.  Their strongest implementation test is exact agreement between the grouped WebAssembly and quantized reference across the retained complete caches and outputs.  Their strongest numerical preservation result is the per-observation common-offset margin check.  The generated texts record repetition and incorrect statements in both models.  The runtime result describes a scalar CPU WebAssembly implementation on one host, with three measured full traces and a stated timing interval.  It supplies no comparison with optimized native matrix libraries, SIMD quantized kernels, GPU inference, or other hardware.

The next numerical task is to tighten the forward bound and discharge captured-operand identity with the source recurrence through checked evidence.  Componentwise bounds, better control of normalization and attention perturbations, and preservation of useful dependencies are possible directions whose benefit requires measurement and proof.  A useful milestone would certify a token from pre-output error bounds and a strict winning margin.  The present experiment establishes the theorem structure and a complete bound evaluation, with zero such certifications.

Repository release completion requires the remaining matcher proof repairs, current aggregate source and frozen-package checks, full execution and semantic-conformance runs, and a receipt from an immutable clean checkout.  The quantized implementation's binary identity and retained GPT evidence are already fixed.  Later release receipts should identify their own source revision and input digest, so a repository-wide result cannot be inferred from the focused package's success.  The report's source and data identifiers permit checking the reported component independently while those broader tasks continue.
